% Generated draft: requires mathematical and source review.
\documentclass[11pt]{article}
\usepackage[margin=25mm]{geometry}
\usepackage{fontspec}
\setmainfont{Latin Modern Roman}
\usepackage{amsmath,amssymb,mathtools,unicode-math}
\setmathfont{Latin Modern Math}
\usepackage{enumitem,needspace,xurl,hyperref}
\setlength{\parindent}{0pt}
\setlength{\parskip}{5pt}
\setlength{\emergencystretch}{3em}
\newcommand{\N}{\mathbb{N}}
\newcommand{\Z}{\mathbb{Z}}
\newcommand{\Q}{\mathbb{Q}}
\newcommand{\R}{\mathbb{R}}
\newcommand{\C}{\mathbb{C}}
\newcommand{\F}{\mathbb{F}}
\newcommand{\A}{\mathbb{A}}
\newcommand{\PP}{\mathbb{P}}
\newcommand{\E}{\mathbb{E}}
\newcommand{\eps}{\varepsilon}
\newcommand{\dd}{\,\mathrm{d}}
\begin{document}
\section*{UnsolvedMath ID: 701157}
\textbf{Source problem code:} catalogue entry
\subsection{Definitions and mathematical statement}
A group \(G\) is word-hyperbolic if it has a finite generating set whose Cayley graph, with its graph metric, has \(\delta\)-thin geodesic triangles for some \(\delta\geq 0\). Here a triangle is \(\delta\)-thin if each side lies in the closed \(\delta\)-neighbourhood of the union of the other two sides. A group has finite exponent if there is an integer \(n\geq 1\) such that every element has \(n\)-th power equal to the identity.

The question proposed by A. Yu. Olshanskii is whether, for every word-hyperbolic group \(G\), there exist a normal subgroup \(N\trianglelefteq G\) and an integer \(n\geq 1\) such that \(N\) is a free group and
\[
\forall g\in G,\qquad g^n\in N.
\]
Equivalently, must every such \(G\) have a free normal subgroup \(N\) for which \(G/N\) has finite exponent? The integer \(n\) may depend on \(G\) and \(N\). No finite-generation or rank restriction on \(N\), and no finiteness requirement on \(G/N\), is stated. We use the standard convention that the trivial group is free of rank zero.
\subsection{Short English statement}
Does every word-hyperbolic group have a free normal subgroup such that all elements of the quotient satisfy one common finite power identity? The supplied registry labels this question open; that label is not an independently verified assessment.
\subsection{Sources}
\begin{itemize}
\item Supplied local UnsolvedMath registry record: permanent ID \(701157\), exact source problem code \texttt{catalogue entry}, title \emph{Kourovka Notebook Problem 15.69}. Upstream revision: \texttt{d371ee917ac9647fbb06cd6fc529deec6f3dda82}.
\item \emph{The Kourovka Notebook (Unsolved Problems in Group Theory)}, 15th issue (2002), Problem 15.69; proposed by A. Yu. Olshanskii. The supplied record also identifies the 21st issue edition (2026), arXiv:1401.0300v48. The collection description names E. I. Khukhro and V. D. Mazurov as editors.
\end{itemize}
\end{document}
